\documentclass[a4paper]{article}
\usepackage{geometry}
\usepackage{graphicx}
\usepackage{natbib}
\usepackage{amsmath}
\usepackage{amssymb}
\usepackage{amsthm}
\usepackage{paralist}
\usepackage{epstopdf}
\usepackage{tabularx}
\usepackage{longtable}
\usepackage{multirow}
\usepackage{multicol}
\usepackage[hidelinks]{hyperref}
\usepackage{fancyvrb}
\usepackage{algorithm}
\usepackage{algorithmic}
\usepackage{float}
\usepackage{paralist}
\usepackage[svgname]{xcolor}
\usepackage{enumerate}
\usepackage{array}
\usepackage{times}
\usepackage{url}
\usepackage{fancyhdr}
\usepackage{comment}
\usepackage{environ}
\usepackage{times}
\usepackage{textcomp}
\usepackage{caption}
\usepackage{color}
\usepackage{xcolor}



\urlstyle{rm}

\setlength\parindent{0pt} % Removes all indentation from paragraphs
\theoremstyle{definition}
\newtheorem{definition}{Definition}[]
\newtheorem{conjecture}{Conjecture}[]
\newtheorem{example}{Example}[]
\newtheorem{theorem}{Theorem}[]
\newtheorem{lemma}{Lemma}
\newtheorem{proposition}{Proposition}
\newtheorem{corollary}{Corollary}

\floatname{algorithm}{Procedure}
\renewcommand{\algorithmicrequire}{\textbf{Input:}}
\renewcommand{\algorithmicensure}{\textbf{Output:}}
\newcommand{\abs}[1]{\lvert#1\rvert}
\newcommand{\norm}[1]{\lVert#1\rVert}
\newcommand{\RR}{\mathbb{R}}
\newcommand{\CC}{\mathbb{C}}
\newcommand{\Nat}{\mathbb{N}}
\newcommand{\br}[1]{\{#1\}}
\DeclareMathOperator*{\argmin}{arg\,min}
\DeclareMathOperator*{\argmax}{arg\,max}
\renewcommand{\qedsymbol}{$\blacksquare$}

\definecolor{dkgreen}{rgb}{0,0.6,0}
\definecolor{gray}{rgb}{0.5,0.5,0.5}
\definecolor{mauve}{rgb}{0.58,0,0.82}

\newcommand{\Var}{\mathrm{Var}}
\newcommand{\Cov}{\mathrm{Cov}}

\newcommand{\vc}[1]{\boldsymbol{#1}}
\newcommand{\xv}{\vc{x}}
\newcommand{\Sigmav}{\vc{\Sigma}}
\newcommand{\alphav}{\vc{\alpha}}
\newcommand{\muv}{\vc{\mu}}

\newcommand{\red}[1]{\textcolor{red}{#1}}

\def\x{\mathbf x}
\def\y{\mathbf y}
\def\w{\mathbf w}
\def\v{\mathbf v}
\def\E{\mathbb E}
\def\V{\mathbb V}

% TO SHOW SOLUTIONS, include following (else comment out):
\newenvironment{soln}{
    \leavevmode\color{blue}\ignorespaces
}{}

\newcommand \mle [1]{{\hat #1}^{\rm MLE}}



\hypersetup{
%    colorlinks,
    linkcolor={red!50!black},
    citecolor={blue!50!black},
    urlcolor={blue!80!black}
}

\geometry{
  top=1in,            % <-- you want to adjust this
  inner=1in,
  outer=1in,
  bottom=1in,
  headheight=3em,       % <-- and this
  headsep=2em,          % <-- and this
  footskip=3em,
}


\pagestyle{fancyplain}
\lhead{\fancyplain{}{Homework 2}}
\rhead{\fancyplain{}{CS 760 Machine Learning}}
\cfoot{\thepage}

\title{\textsc{Homework 2}} % Title

%%% NOTE:  Replace 'NAME HERE' etc., and delete any "\red{}" wrappers (so it won't show up as red)

\author{
\red{$>>$NAME HERE$<<$} \\
\red{$>>$ID HERE$<<$}\\
} 

\date{}

\begin{document}

\maketitle 


\textbf{Instructions:} 

Use this LaTeX file as a template to develop your homework and submit it as a single PDF file on Gradescope. 
For the programming component, you can use any programming language, although the teaching stuff will be most able to assist with Python and specifically the package {\tt scikit-learn}; 
you can use any functionality of the latter.
Put all code used for the assignment at the end of the document in whatever way is easiest but still legible.

\section{Linear Regression and Regularization}
Suppose we model the data-generation process as $y=w^Tx+\epsilon$, where inputs $x$ and an unknown parameter $w$ are $d$-dimensional and $\epsilon\sim\mathcal N(0,\sigma^2$) for some $\sigma\ge0$.

\subsection{Conditional probability [5 points]}
	Write down an expression for $P(y \vert x)$.
	
	{\color{blue}
	Your answer here.
	}
	
\subsection{Conditional log-likelihood [10 points]}
	Assume you are given a set of training observations $(x^{(i)}, y^{(i)})$ for $i = 1,\ldots,n$.   
	Derive the conditional log likelihood of this training data.   
	Drop any constants that do not depend on the parameter $w$.
	
	{\color{blue}
		Your answer here.
	}
	
\subsection{MLE is the least-squares minimizer [5 points]}
	Based on your answer above, show that finding the MLE of that conditional log likelihood is equivalent to minimizing least squares, $\frac{1}{2}\sum\limits_{i=1}^n (y^{(i)} - w^Tx^{(i)})^2 $.
	
	{\color{blue}
			Your answer here.
	}
	
\subsection{Ridge regression optimization [5 points]}
	Suppose we add a Ridge penalty to the objective, i.e. $\frac{1}{2}\sum\limits_{i=1}^n (y^{(i)}-w^Tx^{(i)})^2+\frac\lambda2\|w\|_2^2$ for some $\lambda>0$.
	Although there is a closed form solution to this problem, there are situations in practice where we solve this via gradient descent.
	Suppose at the $t$th step of gradient descent we are at position $w_t\in\mathbb R^d$, and we are using a fixed learning rate $\eta>0$.
	Compute the position $w_{t+1}\in\mathbb{R}^d$ at the $(t+1)$st step.
	
	{\color{blue}
			Your answer here.
	}

\subsection{Effect of the penalty term on optimization [5 points]}

Suppose the data is entirely uninformative, i.e. $x^{(i)}=0_d~\forall~i$, but we still run gradient descent starting from a nonzero initialization. What is the difference in the behavior of gradient descent with and without the Ridge penalty?
Assume $\lambda\le1/\eta$.

{\color{blue}
		Your answer here.
}
	
\subsection{Bayesian view of Ridge regression [15 points]}
	Assume a prior on $w$ that each of its $d$ entries is drawn i.i.d. from $\mathcal N(0,\tau^2)$.
	Prove that the MAP estimate of $w$ with this prior is equivalent to minimizing the above regularized least squares problem with $\lambda = \frac{\sigma^2}{\tau^2}$.  
	Hint: recall that the MAP estimate is the one that maximizes the posterior probability, which is itself the product of the likelihood and the prior probability.
	
	{\color{blue}
			Your answer here.
	}

\newpage
\section{Logistic regression with different representations}

\subsection{Bag-of-words [5 points]}\label{sec:bow}

The files {\tt imdb\_train.csv} and {\tt imdb\_test.csv} both consist of 25K lines, on each of which there is a label ({\tt pos} or {\tt neg}) followed by a tab and followed by a movie review.
We will construct Bag-of-words~(BoW) representations as follows:
\begin{enumerate}
	\item {\bf Tokenize the documents} by lower-casing it, replacing all punctuation (e.g. all characters in Python's {\tt string.punctuation}) with spaces, and replacing document by lists of tokens by splitting along spaces.
	\item {\bf Construct a vocabulary} by taking the 10K tokens that appear most often in the training data lists.
	\item {\bf Represent documents as BoWs} by assigning each token to an index and representing each document in both the training and testing set by a vector $x\in\{0,1\}^{10K}$ s.t. $x_{[j]}=1$ if the document contains the $j$th token and otherwise set $x_{[j]}=0$.
\end{enumerate}

Represent the training data as the pair $(X_\textrm{train},y_\textrm{ytrain})\in\{0,1\}^{25K\times10K}\times\{0,1\}^{25K}$ consisting of the input data $X_\textrm{train}$ and labels $y_\textrm{train}$, and similarly represent the test data as the pair $(X_\textrm{test},y_\textrm{test})\in\{0,1\}^{25K\times10K}\times\{0,1\}^{25K}$.
Train (unregularized)~logistic regression on this training data and report the accuracy of the resulting classifier on the test data.
Hint: if you use {\tt scikit-learn}'s {\tt LogisticRegression}, setting {\tt solver='liblinear'} tends to lead to consistent convergence.

{\color{blue}
		Your answer here.
}

\subsection{Standardization [5 points]}\label{sec:standardization}

Often we standardize the features by applying the transformation $(x_{[j]}-\mu_j)/\sigma_j$ to the $j$th entry of each input $x$, where $\mu_j$ and $\sigma_j$ are the empirical mean and standard deviation of the $j$th feature across the training data.
What is a good reason to avoid doing that here?

{\color{blue}
		Your answer here.
}

\subsection{Regularization and cross-validation [5 points]}

We can do better by regularizing the norm of the linear classifier using the penalty $\tfrac1{Cn}\|w\|_2^2$ on the classification weights $w$, where $C$ is an inverse regularization strength and $n=25K$ is the number of training examples.\footnote{We use inverse regularization strength $C$ instead of the usual regularization strength $\lambda$ because that is how {\tt scikit-learn} does it.}
Train $\ell_2^2$-regularized logistic regression on the training data, with the hyperparameter $C$ selected by 4-fold cross-validation from a set of ten candidates values spaced logarithmically in the interval $[$1E-4,1E4$]$, and report the accuracy of the resulting classifier on the test data.
Hint: you may use {\tt scikit-learn}'s {\tt LogisticRegressionCV}, which defaults to the above candidate values of $C$.

{\color{blue}
		Your answer here.
}

\subsection{Feature selection [5 points]}

Report the value that the linear classifier weight assigns to the tokens ``the'', ``good'', and ``bad'' in both the unregularized and the regularized case.
What does this say about the effect of regularization?

{\color{blue}
		Your answer here.
}

\subsection{Latent Semantic Analysis [5 points]}

The file {\tt imdb\_unsup.txt} contains 50K lines, on each of which there is an unlabeled movie review.
Tokenize it and construct BoWs from it as in Question~\ref{sec:bow} to get a matrix $X\in\{0,1\}^{50K\times10K}$.
To use this data, we will treat tokens as samples and whether or not they appear in document $i$ as features, and obtain a 10-dimensional representation of each token $j$ in the vocabulary by taking the first ten singular values $\Sigma_k$ and singular vectors $U_k$ of $X^T$ and multiplying them together.
This is called Latent Semantic Analysis~(LSA) and is a dimensionality reduction alternative to PCA often used for documents.
Compute $R^{(\textrm{LSA})}=U_k\Sigma_k\in\mathbb{R}^{10K\times 10}$ and {\bf briefly give one practical reason why PCA might not be appropriate to use here}.
Hint: {\tt scikit-learn}'s {\tt TruncatedSVD.fit\_transform} will compute the LSA embeddings directly; if using it, setting {\tt algorithm='arpack'} derandomizes the algorithm and yields more consistent results.

{\color{blue}
	Your answer here.
}

\subsection{Continuous BoW [5 points]}
We now have a low-dimensional representation for each token, but to classify documents with logistic regression we need to represent documents instead.
A simple way to do this is to just represent document $i$ by the sum of the representations $R_{[j]}^{(LSA)}$ of all tokens $j\in[10K]$ that occur in that document.
This is called a Continuous Bag-of-Words (CBoW) representation.
Train and test 4-fold cross-validated logistic regression when the input matrices are the LSA-based CBoW alternatives $X_\textrm{train}^{(LSA)},X_\textrm{tst}^{(LSA)}\in\mathbb{R}^{25K\times10}$ (the label vectors $y_\textrm{train},y_\textrm{test}$ remain the same) and report the test accuracy.

{\color{blue}
	Your answer here.
}

\subsection{Advantages and disadvantages of dimensionality reduction [5 points]}

Ignoring accuracy, list one advantage and one disadvantage of using these LSA-based CBoW representation.

{\color{blue}
	Your answer here.
}

\subsection{Choice of embedding dimension [5 points]}

We specified you use an embedding dimension of 10 for LSA, but list two reasonable ways you could have chosen a suitable dimension using the data.
Hint: one way uses the supervised training data, and another way uses only the unsupervised data.

{\color{blue}
	Your answer here.
}

\subsection{You shall know a word by the company it keeps (Firth, 1957) [5 points]}

The file {\tt glove.csv} contains a token, followed by a tab, followed by a space-delimited list of 300 numbers representing that token.
This is a subsampled set of GloVe word vectors (\url{https://nlp.stanford.edu/projects/glove/}) trained on 840B tokens of Common Crawl web data;
these vectors are produced using a different unsupervised model that counts documents as similar if they co-occur within a small span of words of each other.
Use these GloVe vectors to construct input matrices $X_\textrm{train}^{(GloVe)},X_\textrm{tst}^{(GloVe)}\in\mathbb{R}^{25K\times300}$ consisting of 300-dimensional GloVe-based CBoWs, train and test 4-fold cross-validated logistic regression on them, and report the test accuracy.

{\color{blue}
	Your answer here.
}

\subsection{Learning curves [5 points]}\label{sec:curves}

Train 4-fold cross-validated logistic regression on BoW, LSA-based CBoW, and GloVe-based CBoW document representations while varying the amount of training data provided, specifically using the following number of training examples: 8, 40, 200, 1K, 5K, and 25K.
Make the number of positive and negative samples the same in all datasets and plot the test accuracy vs. number of examples of all three approaches.

{\color{blue}
	Your answer here.
}

\subsection{All you need is GloVe [5 points]}

Discuss what methods perform best at different accuracy levels.
What regimes is unsupervised learning most useful in?

{\color{blue}
	Your answer here.
}

\end{document}
